% ----------------------------------------------------------
\begin{frame}[fragile]{ORM -- Lazy vs Eager Loading}
  When an entity has a relationship to another entity, the ORM must decide
  \emph{when} to load the related object from the database.

  \begin{block}{Two Loading Strategies}
    \begin{itemize}
      \item \textbf{EAGER loading}: the related entity is loaded \emph{immediately} together
            with the parent, via a SQL \texttt{JOIN}. Simple but can load far more data than needed.
            \textbf{Default for \texttt{@ManyToOne} and \texttt{@OneToOne}}.
      \item \textbf{LAZY loading}: the related entity is loaded \emph{on first access} (a second
            query is issued only when \texttt{album.getTracks()} is called).
            \textbf{Default for \texttt{@OneToMany} and \texttt{@ManyToMany}}.
    \end{itemize}
  \end{block}

  \begin{lstlisting}[style=sqlstyle, language=Java]
// Override the default strategy:
@ManyToOne(fetch = FetchType.LAZY)    // LAZY for *ToOne
@JoinColumn(name = "artist_id")
private Artist artist;

@OneToMany(mappedBy = "album", fetch = FetchType.EAGER)
private List<Track> tracks;           // EAGER for *ToMany (rare!)
  \end{lstlisting}
\end{frame}

% ----------------------------------------------------------
\begin{frame}[fragile]{The N+1 Problem -- Explained}
  The N+1 problem is the most common ORM performance pitfall.

  \begin{lstlisting}[style=sqlstyle, language=Java]
// Suppose we load all albums (1 query):
List<Album> albums = albumRepo.findAll();   // SELECT * FROM album

// Then we iterate and access each album's artist (LAZY):
for (Album album : albums) {
    System.out.println(album.getArtist().getName());
    // Each call fires: SELECT * FROM artist WHERE id = ?
}
// With 100 albums -> 1 + 100 = 101 queries!  (N+1 problem)
  \end{lstlisting}

  \begin{alertblock}{Why Is This a Problem?}
    With 1000 albums, you execute 1001 database queries instead of 1 or 2.
    Each query has network latency, parse overhead, and execution cost.
    In production with millions of records, this can bring an application to its knees.
  \end{alertblock}
\end{frame}

% ----------------------------------------------------------
\begin{frame}[fragile]{Fixing the N+1 Problem -- JOIN FETCH}
  \begin{lstlisting}[style=sqlstyle, language=Java]
// Fix 1: JOIN FETCH in JPQL -- loads everything in ONE query
public interface AlbumRepository extends JpaRepository<Album, Long> {

    @Query("SELECT al FROM Album al JOIN FETCH al.artist")
    List<Album> findAllWithArtist();
    // Generates: SELECT al.*, ar.* FROM album al
    //            JOIN artist ar ON al.artist_id = ar.id
}

// Fix 2: EntityGraph -- declarative eager loading for a query
@EntityGraph(attributePaths = {"artist", "tracks"})
List<Album> findByYear(int year);
// Loads albums + their artists + their tracks in one shot
  \end{lstlisting}

  \begin{block}{Rule of Thumb}
    \begin{itemize}
      \item Default: use \texttt{LAZY} on all relationships.
      \item For specific queries where you \emph{know} you need related data: use
            \texttt{JOIN FETCH} or \texttt{@EntityGraph}.
      \item \textbf{Never} change a relationship to \texttt{EAGER} globally just to fix an N+1
            -- it will cause over-fetching everywhere else.
    \end{itemize}
  \end{block}
\end{frame}

% ----------------------------------------------------------
\begin{frame}{Loading Strategies -- Comparison}
  \begin{center}
  \renewcommand{\arraystretch}{1.35}
  \begin{tabularx}{\linewidth}{@{}lXX@{}}
    \toprule
    \textbf{Strategy}  & \textbf{When data is loaded} & \textbf{Risk} \\
    \midrule
    EAGER (default @ManyToOne)
      & Always, immediately with the parent -- even if never used
      & Over-fetching; cartesian product explosion with multiple EAGER collections \\[4pt]
    LAZY (default @OneToMany)
      & On first access via getter
      & N+1 queries if accessed in a loop; \texttt{LazyInitializationException} outside session \\[4pt]
    JOIN FETCH (query-level)
      & Only for this specific query
      & Best practice; no global side effects \\[4pt]
    @EntityGraph
      & Declared per query/method
      & Best practice; declarative and reusable \\
    \bottomrule
  \end{tabularx}
  \end{center}
\end{frame}
